\begin{lemma}
For every $\eta>0$ there are $\Delta_0$ and $\gamma_0$ such that the
following holds.  Let $G$ be a $\Delta$-regular graph with
$\Delta\ge\Delta_0$, let $\gamma\ge\gamma_0$, and let $I$ be sampled by
\noderef{RandomIndependentSetSampling}.  Fix $r$ and
$X\subseteq N_G(r)$, and assume that no two distinct vertices of $X$ have a
common neighbor outside $N_G[r]$.  Let $\mathcal J_3$ be the independent
triples in $X$, and, for every independent pair $ab\subseteq X$, put
\[
\ell_{ab}=\frac{|N_G(a)\cap N_G(b)|}{\Delta}.
\]
Then
\[
\mathbb E\binom{|I\cap X|}{3}
\le
\frac{|\mathcal J_3|}{\Delta^3}
+\frac{1}{3\Delta^3}
\sum_{abc\in\mathcal J_3}\ \sum_{\{p,q\}\subset\{a,b,c\}}
\left(\frac{2}{(2-\ell_{pq})(1-\ell_{pq})}-1\right)
+\eta .
\]
\end{lemma}
\begin{proof}
Write $\overline T_X=\binom{|I\cap X|}{3}$.  By linearity of expectation,
\[
\mathbb E\overline T_X
=\sum_{Q\in\mathcal J_3}\Pr[Q\subseteq I],
\]
because the output set in \noderef{RandomIndependentSetSampling} is always
independent, so a three-subset of $X$ can contribute only if it is one of the
triples in $\mathcal J_3$.

Fix one triple $Q=\{a,b,c\}\in\mathcal J_3$.  Condition on the event that
$a,b,c$ are activated, and use the rescaled priority variables
$x_v=\gamma(1-\pi(v))$ for $v\in Q$.  The three activation probabilities and
the uniform priority law in \noderef{RandomIndependentSetSampling} give the
factor $dx_a\,dx_b\,dx_c/\Delta^3$ on $[0,\gamma]^3$.  The sampling
definition also gives $0<\gamma$ and $0\le\gamma/\Delta\le1$, because
$\gamma/\Delta$ is the Bernoulli activation probability.  Hence
$0<\gamma\le\Delta$, and every quantity $x_v/\Delta$ used below lies in
$[0,1]$.  A vertex
$z\notin Q$ adjacent to the nonempty set $S_z\subseteq Q$ deletes a member of
$Q$ exactly when $z$ is activated and has priority larger than at least one
neighbor in $S_z$, and this has conditional probability
$\max_{v\in S_z}x_v/\Delta$.  Hence the conditional survival probability is
the product over such vertices of
$1-\max_{v\in S_z}x_v/\Delta$.

We now estimate this product on each ordered chamber.  On the chamber
$x_a\ge x_b\ge x_c$, classify vertices outside $Q$ according to their
nonempty adjacency pattern to $Q$:
\[
n_a,n_b,n_c,n_{ab},n_{ac},n_{bc},n_{abc}.
\]
Thus, for example, $n_{ab}$ counts vertices adjacent to $a$ and $b$ but not
to $c$.  The clean-neighborhood hypothesis says that for every pair in $Q$
all common neighbors are in $N_G[r]$; in particular
\[
\ell_{ab}=\frac{n_{ab}+n_{abc}}{\Delta},\quad
\ell_{ac}=\frac{n_{ac}+n_{abc}}{\Delta},\quad
\ell_{bc}=\frac{n_{bc}+n_{abc}}{\Delta}.
\]
The $\Delta$ neighbors of $a$ all contribute factors with maximum coordinate
$x_a$.  After those have been counted, the remaining neighbors of $b$ are
exactly the $b$-only and $bc$-vertices, namely
$\Delta-(n_{ab}+n_{abc})=(1-\ell_{ab})\Delta$ vertices, and they contribute
with maximum coordinate $x_b$.  Finally, after all vertices adjacent to
$a$ or $b$ have been counted, the only remaining neighbors of $c$ are the
$c$-only vertices, namely
\[
n_c=\Delta-n_{ac}-n_{bc}-n_{abc}
=\Delta(1-\ell_{ac}-\ell_{bc})+n_{abc}.
\]
Using $1-t\le e^{-t}$ for $0\le t\le1$, the chamber contribution is at most
\[
\int_{0\le z\le y\le x\le\gamma}
e^{-x}e^{-(1-\ell_{ab})y}e^{-(n_c/\Delta)z}\,dz\,dy\,dx .
\]
The five other chambers have the identical form with the ordered triple
permuted.

It remains to bound the sum of these six elementary integrals.  Set
\[
\lambda_{ab}=1-\ell_{ab},\qquad
\rho_c=n_c/\Delta
=1-\ell_{ac}-\ell_{bc}+n_{abc}/\Delta,
\]
and define $\rho_a,\rho_b$ by cyclic permutation.  The pattern counts are
nonnegative.  For any pair from $Q$, the clean-neighborhood hypothesis puts
all its common neighbors in $N_G[r]$, and independence of $Q$ excludes the
other two vertices of the pair and the third vertex of $Q$ from being common
neighbors.  Hence the pair has at most $\Delta-1$ common neighbors.  Thus
$0<\lambda_{pq}\le1$, while
$0\le\rho_v\le1$ is immediate from its definition as a normalized count.
For fixed larger pair $ab$ the two chambers in which $a$ and $b$ are the two
larger coordinates contribute at most
\[
C_{ab}
=2\int_{0\le z\le y\le x<\infty}
e^{-x}e^{-\lambda_{ab}y}e^{-\rho_c z}\,dz\,dy\,dx .
\]
Direct integration gives, with the value at $\rho_c=0$ understood as the
limit,
\[
\begin{aligned}
C_{ab}
&=2\int_0^\infty e^{-(1+\lambda_{ab})y}
  \left(\int_0^y e^{-\rho_c z}\,dz\right)\,dy  \\
&=\frac{2}{\rho_c}
  \left(\frac1{1+\lambda_{ab}}
  -\frac1{1+\lambda_{ab}+\rho_c}\right)  \\
&=\frac{2}{(1+\lambda_{ab})(1+\lambda_{ab}+\rho_c)}.
\end{aligned}
\]
The cyclic formula gives $C_{ac}$ and $C_{bc}$.

Now put
\[
x=\lambda_{ab},\qquad y=\lambda_{ac},\qquad z=\lambda_{bc},\qquad
\tau=n_{abc}/\Delta .
\]
The degree equations give the common-denominator identities
\[
1+\lambda_{ab}+\rho_c
=1+\lambda_{ac}+\rho_b
=1+\lambda_{bc}+\rho_a
=x+y+z+\tau .
\]
For instance,
\[
1+\lambda_{ab}+\rho_c
=1+\left(1-\frac{n_{ab}+n_{abc}}{\Delta}\right)
  +\left(1-\frac{n_{ac}+n_{bc}+n_{abc}}{\Delta}\right)
=3-\frac{n_{ab}+n_{ac}+n_{bc}+2n_{abc}}{\Delta},
\]
and the two other expressions are the same.  Therefore, writing
$M=x+y+z+\tau$, the sum of the six chamber bounds is
\[
C_{ab}+C_{ac}+C_{bc}
=\frac2M\left(\frac1{1+x}+\frac1{1+y}+\frac1{1+z}\right).
\]
Since $\tau\ge0$, we have $M\ge x+y+z$.  Hence
\[
C_{ab}+C_{ac}+C_{bc}
\le
\frac2{x+y+z}\left(\frac1{1+x}+\frac1{1+y}+\frac1{1+z}\right).
\]
It remains only to compare this last expression with the desired average.
For positive $x,y,z$ the following identity is obtained by putting the left
side over the common denominator
$xy z(1+x)(1+y)(1+z)$ and collecting the three unordered pairs:
\[
\begin{aligned}
&(x+y+z)\left(\frac1{x(1+x)}+\frac1{y(1+y)}
  +\frac1{z(1+z)}\right)
-3\left(\frac1{1+x}+\frac1{1+y}+\frac1{1+z}\right)\\
&\quad =
\frac{(x-y)^2(x+y+1)}{xy(1+x)(1+y)}
+\frac{(x-z)^2(x+z+1)}{xz(1+x)(1+z)}
+\frac{(y-z)^2(y+z+1)}{yz(1+y)(1+z)} .
\end{aligned}
\]
Each term on the right is nonnegative.  Thus
\[
\frac2{x+y+z}\left(\frac1{1+x}+\frac1{1+y}+\frac1{1+z}\right)
\le
\frac23\left(\frac1{x(1+x)}+\frac1{y(1+y)}+\frac1{z(1+z)}\right).
\]
Returning to $\lambda_{pq}=1-\ell_{pq}$, this is exactly
\[
C_{ab}+C_{ac}+C_{bc}
\le
1+\frac13\sum_{pq\in\{ab,ac,bc\}}
\left(\frac{2}{(1+\lambda_{pq})\lambda_{pq}}-1\right).
\]
Indeed the right hand side simplifies to
$\frac23\sum_{pq}1/((1+\lambda_{pq})\lambda_{pq})$ because
$1+\frac13\sum(A_{pq}-1)=\frac13\sum A_{pq}$.
This is precisely where the three ordered chamber families are averaged; no
common-neighbor factor is reused to create an extra exponent.

Combining the six ordered-chamber bounds and restricting from the infinite
orthant back to $[0,\gamma]^3$ gives
\[
\Pr[Q\subseteq I]
\le
\frac1{\Delta^3}
\left[
1+\frac13\sum_{\{p,q\}\subset Q}
\left(\frac{2}{(2-\ell_{pq})(1-\ell_{pq})}-1\right)
\right].
\]
There is no truncation error here: the exponential product is an upper bound
on the original finite product on the whole valid box $[0,\gamma]^3$, and
replacing that box by the infinite chamber only increases the integral.
Summing the last displayed estimate over $Q\in\mathcal J_3$ proves the same
bound without the final $+\eta$ term.  The stated version with $+\eta$ follows
immediately for the given $\eta>0$; for example, take
$\Delta_0=\gamma_0=1$.
\end{proof}
